% !TEX program = pdflatex
\documentclass[11pt,a4paper]{article}

% ─── Geometry ────────────────────────────────────────────────
\usepackage[margin=2.5cm]{geometry}

% ─── Encoding & fonts ───────────────────────────────────────
\usepackage[T1]{fontenc}
\usepackage[utf8]{inputenc}
\usepackage{lmodern}
\usepackage{microtype}

% ─── Math ────────────────────────────────────────────────────
\usepackage{amsmath,amssymb,amsthm}

% ─── Tables & figures ───────────────────────────────────────
\usepackage{booktabs}
\usepackage{tabularx}
\usepackage{multirow}
\usepackage{graphicx}
\usepackage{float}
\usepackage{caption}
\usepackage{subcaption}

% ─── Algorithms ──────────────────────────────────────────────
\usepackage[ruled,linesnumbered,noend]{algorithm2e}
\SetKwComment{Comment}{$\triangleright$~}{}

% ─── Hyperref ────────────────────────────────────────────────
\usepackage[colorlinks,citecolor=blue!70!black,linkcolor=blue!70!black,urlcolor=blue!50!black]{hyperref}
\usepackage[capitalise,noabbrev]{cleveref}

% ─── Code listings ───────────────────────────────────────────
\usepackage{listings}
\lstset{
  basicstyle=\ttfamily\small,
  breaklines=true,
  frame=single,
  columns=fullflexible,
}

% ─── Theorem environments ────────────────────────────────────
\newtheorem{theorem}{Theorem}[section]
\newtheorem{lemma}[theorem]{Lemma}
\newtheorem{corollary}[theorem]{Corollary}
\newtheorem{proposition}[theorem]{Proposition}
\theoremstyle{definition}
\newtheorem{definition}[theorem]{Definition}
\newtheorem{remark}[theorem]{Remark}

% ─── Notation macros ────────────────────────────────────────
\newcommand{\secpar}{\lambda}
\newcommand{\negl}{\mathsf{negl}}
\newcommand{\ppt}{\mathsf{PPT}}
\newcommand{\adv}{\mathbf{Adv}}
\newcommand{\advA}{\mathcal{A}}
\newcommand{\advB}{\mathcal{B}}
\newcommand{\advS}{\mathcal{S}}
\newcommand{\challenger}{\mathcal{C}}
\newcommand{\oracle}{\mathcal{O}}
\newcommand{\HKDF}{\textsf{HKDF}}
\newcommand{\HMAC}{\textsf{HMAC}}
\newcommand{\SHA}{\textsf{SHA-256}}
\newcommand{\AEAD}{\textsf{AES-256-GCM-SIV}}
\newcommand{\DH}{\textsf{X25519}}
\newcommand{\KEM}{\textsf{Kyber-768}}
\newcommand{\KeyGen}{\textsf{KeyGen}}
\newcommand{\Encaps}{\textsf{Encaps}}
\newcommand{\Decaps}{\textsf{Decaps}}
\newcommand{\Enc}{\textsf{Enc}}
\newcommand{\Dec}{\textsf{Dec}}
\newcommand{\GapCDH}{\textsf{Gap-CDH}}
\newcommand{\INDCCA}{\textsf{IND-CCA2}}
\newcommand{\INDCPA}{\textsf{IND-CPA}}
\newcommand{\INTCTXT}{\textsf{INT-CTXT}}
\newcommand{\MRAE}{\textsf{MRAE}}
\newcommand{\PRF}{\textsf{PRF}}
\newcommand{\dualPRF}{\textsf{dual-PRF}}
\newcommand{\ROM}{\textsf{ROM}}
\newcommand{\eCK}{\textsf{eCK}}
\newcommand{\concat}{\,\|\,}
\newcommand{\getsr}{\stackrel{\$}{\gets}}
\newcommand{\rk}{\mathit{rk}}
\newcommand{\ck}{\mathit{ck}}
\newcommand{\mk}{\mathit{mk}}
\newcommand{\mdk}{\mathit{mdk}}
\newcommand{\epoch}{\mathit{epoch}}
\newcommand{\IK}{\mathit{IK}}
\newcommand{\SPK}{\mathit{SPK}}
\newcommand{\EK}{\mathit{EK}}
\newcommand{\OPK}{\mathit{OPK}}

% ─── Title ───────────────────────────────────────────────────
\title{\textbf{Ecliptix: A Hybrid Post-Quantum Secure Messaging Protocol\\with Per-Ratchet KEM Rotation and Formal Verification}}

\author{%
  Oleksandr Melnychenko\\
  \texttt{oleksandr@ecliptix.io}
}

\date{\today}

% ═════════════════════════════════════════════════════════════
\begin{document}
\maketitle

% ─────────────────────────────────────────────────────────────
% ABSTRACT
% ─────────────────────────────────────────────────────────────
\begin{abstract}
We present \textsc{Ecliptix}, a hybrid post-quantum secure messaging protocol that combines classical Diffie--Hellman key agreement (X25519) with lattice-based key encapsulation (Kyber-768/ML-KEM) in both the initial handshake and every subsequent ratchet step. Building on the Signal protocol's Double Ratchet architecture, Ecliptix introduces three principal enhancements: (1)~per-direction hybrid ratcheting, where every direction change triggers a fresh X25519~+~Kyber-768 key exchange, providing post-compromise security that heals in one classical step or two hybrid steps; (2)~a two-layer AEAD construction using AES-256-GCM-SIV with independent per-epoch metadata encryption keys, providing nonce-misuse resistance and metadata privacy; and (3)~augmented HKDF-based key derivation that cryptographically binds each new Kyber public key into the ratchet chain, providing implicit KEM authentication without additional signatures.

We prove six security theorems via game-based reductions under standard assumptions (Gap-CDH, ML-KEM IND-CCA2, dual-PRF, MRAE), establishing eCK-AKE security, forward secrecy, post-compromise security, message confidentiality with integrity, and replay resistance. We complement these proofs with machine-checked formal verification: 10 out of 10 Tamarin Prover lemmas and 4 out of 6 ProVerif queries are verified, with the two unproven ProVerif queries attributed to known DH equational theory overapproximation limitations.

Our Rust implementation comprises 5{,}400 lines of code, passes 159 tests (including property-based and concurrent stress tests), achieves sub-millisecond handshake latency ($\sim$1.1\,ms) and $\sim$15\,$\mu$s per-message throughput on commodity hardware, and exports a C~FFI for cross-platform deployment.
\end{abstract}

\smallskip
\noindent\textbf{Keywords:} Post-quantum cryptography, secure messaging, double ratchet, hybrid key exchange, Kyber, formal verification, Tamarin, ProVerif.

\tableofcontents
\newpage

% ═════════════════════════════════════════════════════════════
% 1. INTRODUCTION
% ═════════════════════════════════════════════════════════════
\section{Introduction}\label{sec:intro}

The Signal protocol~\cite{signal-x3dh,cohn-gordon2020} has become the de facto standard for end-to-end encrypted messaging, underpinning WhatsApp, Signal Messenger, and Google Messages. Its core innovations---the X3DH key agreement and the Double Ratchet algorithm---provide strong forward secrecy and post-compromise security guarantees using elliptic-curve Diffie--Hellman key exchange.

However, the advent of cryptographically relevant quantum computers (CRQCs) poses an existential threat to protocols relying exclusively on classical assumptions. While NIST has standardized the ML-KEM (Kyber) lattice-based key encapsulation mechanism~\cite{fips203}, integrating post-quantum primitives into messaging protocols presents several challenges:

\paragraph{Challenge 1: Harvest-now-decrypt-later attacks.}
An adversary recording encrypted traffic today can decrypt it once a CRQC becomes available. Forward secrecy based solely on ephemeral ECDH provides no protection against this threat, as Shor's algorithm efficiently solves the discrete logarithm problem on elliptic curves.

\paragraph{Challenge 2: KEM integration granularity.}
The Signal PQXDH proposal~\cite{signal-pqxdh} introduces a single Kyber encapsulation during the initial handshake but does not include post-quantum key material in subsequent ratchet steps. This means that post-compromise security recovery relies entirely on classical DH, offering no quantum resistance after the initial key agreement.

\paragraph{Challenge 3: Asymmetric KEM recovery.}
Unlike DH, where both parties contribute ephemeral keys symmetrically, KEM requires one party to advertise a public key before the other can encapsulate. This introduces an inherent asymmetry in hybrid ratchet protocols: classical PCS recovery occurs in one ratchet step, but full hybrid recovery requires two steps.

\paragraph{Challenge 4: Metadata leakage.}
Standard Double Ratchet implementations transmit ratchet public keys, epoch counters, and message indices in cleartext envelope headers. Even with payload encryption, this metadata reveals communication patterns, message ordering, and session lifecycle to network observers.

\medskip
\noindent\textbf{Our contributions.} We present \textsc{Ecliptix}, a hybrid post-quantum secure messaging protocol that addresses all four challenges:

\begin{enumerate}
  \item \textbf{Per-direction hybrid ratcheting.} Every direction change triggers a full X25519~+~Kyber-768 key exchange, generating fresh classical and post-quantum key material. This provides post-compromise security that heals in one step for classical adversaries and two steps for quantum adversaries (\cref{thm:pcs}).

  \item \textbf{Augmented HKDF key derivation.} We bind each new Kyber public key into the HKDF info parameter during ratchet key derivation, providing implicit authentication of the KEM public key without requiring additional digital signatures (\cref{lem:augmented-hkdf}).

  \item \textbf{Two-layer AEAD with metadata encryption.} We employ AES-256-GCM-SIV---a nonce-misuse-resistant authenticated encryption scheme---in a two-layer construction where envelope metadata is encrypted with an independent per-epoch key, providing metadata privacy with defense-in-depth against nonce reuse (\cref{lem:two-layer}).

  \item \textbf{Comprehensive formal verification.} We provide game-based security proofs (6 theorems, 8 lemmas), machine-checked Tamarin Prover models (10/10 lemmas verified), and ProVerif analysis (4/6 queries proven), offering the strongest formal assurance of any hybrid post-quantum messaging protocol to date.

  \item \textbf{Production-ready implementation.} Our Rust implementation features secure memory management via \texttt{mlock} and \texttt{zeroize}, 337+ automated tests including property-based testing, and a C~FFI layer for cross-platform deployment.
\end{enumerate}

\paragraph{Paper organization.}
\Cref{sec:related} surveys related work. \Cref{sec:prelim} establishes notation and cryptographic definitions. \Cref{sec:protocol} describes the Ecliptix protocol in detail. \Cref{sec:security-model} defines the security model. \Cref{sec:proofs} presents the six main theorems with proof sketches. \Cref{sec:formal} describes the formal verification methodology and results. \Cref{sec:implementation} details the implementation. \Cref{sec:evaluation} presents performance benchmarks and test coverage. \Cref{sec:comparison} compares Ecliptix with Signal and PQXDH. \Cref{sec:discussion} discusses limitations and future work. \Cref{sec:conclusion} concludes.


% ═════════════════════════════════════════════════════════════
% 2. RELATED WORK
% ═════════════════════════════════════════════════════════════
\section{Related Work}\label{sec:related}

\paragraph{Signal protocol.}
The Signal protocol, formalized by Cohn-Gordon et al.~\cite{cohn-gordon2020}, achieves forward secrecy and post-compromise security through the combination of X3DH key agreement~\cite{signal-x3dh} and the Double Ratchet algorithm. The protocol has received extensive formal analysis in the symbolic~\cite{cohn-gordon2020} and computational~\cite{lamacchia2007} models. However, Signal relies exclusively on classical Diffie--Hellman assumptions (Curve25519), offering no protection against quantum adversaries.

\paragraph{Post-quantum extensions to Signal.}
Signal's PQXDH protocol~\cite{signal-pqxdh} adds a single Kyber-768 encapsulation to the X3DH handshake, providing initial quantum resistance. Brendel et al.~\cite{brendel2020} propose a generic post-quantum X3DH construction, and Hashimoto et al.~\cite{hashimoto2021} provide an efficient instantiation. These works focus on the handshake phase; none integrate post-quantum key material into the ratchet mechanism.

\paragraph{Post-quantum ratcheting.}
Alwen et al.~\cite{alwen2020} and Bienstock et al.~\cite{bienstock2022} study post-quantum ratcheting in the context of continuous group key agreement (CGKA), focusing on tree-based constructions for group messaging (MLS). Their security models differ significantly from the two-party setting we consider.

\paragraph{Hybrid key exchange.}
The concept of hybrid key exchange---combining classical and post-quantum primitives---has been studied in the TLS context~\cite{stebila2020,dowling2020}. These works establish security under the assumption that at least one component remains secure. Our hybrid combiner (\cref{thm:hybrid-combiner}) follows this approach, using HKDF with the classical shared secret as salt and the KEM shared secret as input keying material.

\paragraph{Formal verification of messaging protocols.}
Tamarin~\cite{meier2013} and ProVerif~\cite{blanchet2016} have been used to verify properties of the Signal protocol~\cite{cohn-gordon2020}. Our work extends these analyses to the hybrid post-quantum setting, introducing custom equational theories for Kyber-768 and addressing the Tamarin-specific challenges of source saturation with multiple DH operations.

\paragraph{Nonce-misuse-resistant AEAD.}
AES-256-GCM-SIV~\cite{rfc8452} provides misuse-resistant authenticated encryption, degrading gracefully to deterministic encryption under nonce reuse rather than catastrophically leaking plaintext as in standard AES-GCM. Rogaway and Shrimpton~\cite{rogaway2006} formalize the MRAE security notion. We employ GCM-SIV as defense-in-depth alongside monotonic nonce counters.


% ═════════════════════════════════════════════════════════════
% 3. PRELIMINARIES
% ═════════════════════════════════════════════════════════════
\section{Preliminaries}\label{sec:prelim}

\subsection{Notation}

We write $\secpar$ for the security parameter. A function $f(\secpar)$ is \emph{negligible} if $f(\secpar) \in o(\secpar^{-c})$ for every constant $c > 0$; we write $f \in \negl(\secpar)$. All algorithms are implicitly parameterized by $\secpar$. We write $x \getsr S$ for uniform random sampling from a finite set $S$. The operator $\concat$ denotes byte-string concatenation. We use $\ppt$ to denote probabilistic polynomial-time.

\begin{table}[h]
\centering
\caption{Protocol-specific notation.}\label{tab:notation}
\small
\begin{tabular}{@{}ll@{}}
\toprule
Symbol & Meaning \\
\midrule
$\IK_X$ & Long-term X25519 identity keypair of party $X$ \\
$\SPK_X$ & Signed pre-key (X25519) of party $X$ \\
$\EK_X$ & Ephemeral X25519 keypair of party $X$ \\
$\OPK_X$ & One-time pre-key (X25519) of party $X$ \\
$(kpk_X, ksk_X)$ & Kyber-768 keypair of party $X$ \\
$\rk_e$ & Root key at epoch $e$ \\
$\ck_e$ & Chain key at epoch $e$ \\
$\mk_{e,n}$ & Message key at epoch $e$, chain index $n$ \\
$\mdk_e$ & Metadata key at epoch $e$ \\
$R(\cdot)$ & Hybrid ratchet step function \\
$q_s$ & Number of sessions \\
$q_m$ & Number of messages per session \\
$N_{\max}$ & Maximum chain length (default 10{,}000) \\
\bottomrule
\end{tabular}
\end{table}

\subsection{Cryptographic Primitives}

\begin{definition}[Diffie--Hellman on Curve25519]\label{def:dh}
We work over the Montgomery curve Curve25519 of prime order $q = 2^{252} + 27742317777372353535851937790883648493$. The X25519 function~\cite{rfc7748} computes scalar multiplication $\DH(a, B) = [a]B$. A keypair is $(a, A)$ where $A = [a]G$ and $G$ is the base point $u = 9$.
\end{definition}

\begin{definition}[Gap-CDH]\label{def:gap-cdh}
The $\GapCDH$ assumption states that for all $\ppt$ adversaries $\advA$ with access to a DDH oracle $\oracle_{\mathrm{DDH}}$:
\[
  \adv^{\GapCDH}_G(\advA) = \Pr\!\left[\advA^{\oracle_{\mathrm{DDH}}}(G, [a]G, [b]G) = [ab]G \;\middle|\; a, b \getsr \mathbb{Z}_q\right] \leq \negl(\secpar).
\]
\end{definition}

\begin{definition}[Key Encapsulation Mechanism]\label{def:kem}
A KEM $\Pi = (\KeyGen, \Encaps, \Decaps)$ consists of key generation outputting $(pk, sk)$, encapsulation $\Encaps(pk) \to (ct, ss)$, and decapsulation $\Decaps(sk, ct) \to ss$, satisfying correctness: $\Decaps(sk, ct) = ss$ for $(ct, ss) \gets \Encaps(pk)$.
\end{definition}

\begin{definition}[IND-CCA2 for KEM]\label{def:ind-cca2}
The $\INDCCA$ advantage of $\advA$ against KEM $\Pi$ is:
\[
  \adv^{\INDCCA}_\Pi(\advA) = \left|\Pr\!\left[\advA^{\oracle_{\Decaps}}(pk, ct^*, ss_b) = b\right] - \frac{1}{2}\right| \leq \negl(\secpar),
\]
where $b \getsr \{0,1\}$, $(ct^*, ss_0) \gets \Encaps(pk)$, $ss_1 \getsr \{0,1\}^{256}$, and $\advA$ may not query $\oracle_{\Decaps}$ on $ct^*$.
\end{definition}

\begin{definition}[HKDF Dual-PRF Property~\cite{krawczyk2010}]\label{def:dual-prf}
$\HKDF$-$\SHA$ satisfies the $\dualPRF$ property if:
\begin{enumerate}
  \item \textbf{PRF in key (salt):} For a uniformly random salt $s$, $\HKDF(s, \cdot)$ is computationally indistinguishable from a random function.
  \item \textbf{PRF in input (IKM):} For a uniformly random IKM $k$, $\HKDF(\cdot, k)$ is computationally indistinguishable from a random function.
\end{enumerate}
\end{definition}

\begin{definition}[MRAE Security~\cite{rfc8452,rogaway2006}]\label{def:mrae}
$\AEAD$ satisfies $\MRAE$ security if:
\begin{itemize}
  \item Under unique nonces: full $\INDCPA$ + $\INTCTXT$ (standard AEAD security).
  \item Under nonce reuse: degrades to deterministic authenticated encryption (DAE)---only the equality of plaintexts encrypted under the same nonce leaks; $\INTCTXT$ is preserved.
\end{itemize}
\end{definition}


% ═════════════════════════════════════════════════════════════
% 4. PROTOCOL DESIGN
% ═════════════════════════════════════════════════════════════
\section{Protocol Design}\label{sec:protocol}

Ecliptix consists of two phases: a \emph{hybrid X3DH handshake} (\cref{sec:handshake}) that establishes initial shared state, and a \emph{hybrid Double Ratchet} (\cref{sec:ratchet}) that continuously refreshes key material. Messages are protected by a \emph{two-layer AEAD} construction (\cref{sec:two-layer}).

\begin{figure}[t]
\centering
\begin{minipage}{0.95\textwidth}
\begin{verbatim}
  Handshake (Hybrid X3DH)          Session (Hybrid Double Ratchet)
  ┌─────────────────────────┐      ┌──────────────────────────────────┐
  │ 4× X25519 DH            │      │ Per-direction ratchet:           │
  │ 1× Kyber-768 KEM        │─────>│   X25519 DH + Kyber-768 KEM     │
  │ HKDF-SHA256 combiner    │      │ Chain KDF: HKDF-SHA256           │
  │ HMAC key confirmation   │      │ AEAD: AES-256-GCM-SIV           │
  │ Ed25519 SPK signature   │      │ Metadata: independent AEAD layer │
  └─────────────────────────┘      └──────────────────────────────────┘
\end{verbatim}
\end{minipage}
\caption{High-level architecture of the Ecliptix protocol.}\label{fig:architecture}
\end{figure}

\subsection{Pre-Key Bundle}\label{sec:bundle}

Each party $R$ publishes a pre-key bundle:
\[
  \mathit{Bundle}_R = \bigl(\IK_R^{\mathit{pub}},\; \SPK_R^{\mathit{pub}},\; \sigma_{\SPK} = \textsf{Ed25519.Sign}(ik_R^{\mathit{ed}}, \SPK_R^{\mathit{pub}}),\; \{\OPK_{R,j}^{\mathit{pub}}\},\; kpk_R\bigr),
\]
where $\IK_R$ is the long-term X25519 identity key, $\SPK_R$ is a medium-term signed pre-key (signed with the corresponding Ed25519 identity key), $\{\OPK_{R,j}\}$ is a set of one-time pre-keys, and $kpk_R$ is the Kyber-768 public key.

\subsection{Hybrid X3DH Key Agreement}\label{sec:handshake}

The handshake extends Signal's X3DH with a Kyber-768 KEM encapsulation and bilateral key confirmation.

\begin{algorithm}[H]
\caption{Initiator X3DH computation.}\label{alg:initiator}
\DontPrintSemicolon
\KwIn{Initiator identity $(\IK_I, \EK_I)$; responder bundle $\mathit{Bundle}_R$}
\KwOut{Root key $\rk_0$, session state, init message}

Validate $\mathit{Bundle}_R$: verify $\sigma_{\SPK}$, reject small-order points, validate $kpk_R$\;
Generate ephemeral keypair $(\mathit{ek}_I, \EK_I^{\mathit{pub}})$\;
Select one-time pre-key $\OPK_R$ (if available)\;
\BlankLine
\Comment{Four X25519 DH computations}
$\mathit{dh}_1 \gets \DH(\mathit{ik}_I, \SPK_R^{\mathit{pub}})$\;
$\mathit{dh}_2 \gets \DH(\mathit{ek}_I, \IK_R^{\mathit{pub}})$\;
$\mathit{dh}_3 \gets \DH(\mathit{ek}_I, \SPK_R^{\mathit{pub}})$\;
$\mathit{dh}_4 \gets \DH(\mathit{ek}_I, \OPK_R^{\mathit{pub}})$ \hfill\Comment{if OPK present}
\BlankLine
\Comment{Classical shared secret}
$\mathit{IKM} \gets \mathtt{0xFF}^{32} \concat \mathit{dh}_1 \concat \mathit{dh}_2 \concat \mathit{dh}_3 [\concat \mathit{dh}_4]$\;
$\mathit{cs} \gets \HKDF(\mathit{IKM},\, 32,\, \varepsilon,\, \texttt{"Ecliptix-X3DH"})$\;
\BlankLine
\Comment{Kyber-768 encapsulation}
$(\mathit{ct}_{\mathit{ky}},\, \mathit{ss}_{\mathit{ky}}) \gets \KEM.\Encaps(kpk_R)$\;
\BlankLine
\Comment{Hybrid combiner (classical as salt, PQ as IKM)}
$\rk_0 \gets \HKDF(\mathit{ss}_{\mathit{ky}},\, 32,\, \mathit{cs},\, \texttt{"Ecliptix-Hybrid-X3DH"})$\;
\BlankLine
\Comment{Derive session keys}
$\mathit{sid} \gets \HKDF(\rk_0,\, 16,\, \varepsilon,\, \texttt{"Ecliptix-SessionId"})$\;
$\mdk_0 \gets \HKDF(\rk_0,\, 32,\, \mathit{ctx}_{\mathit{md}},\, \texttt{"Ecliptix-MetadataKey"})$\;
$\mathit{kc}_I \gets \HKDF(\rk_0,\, 32,\, \varepsilon,\, \texttt{"Ecliptix-KeyConfirm-I"})$\;
$\mathit{kc}_R \gets \HKDF(\rk_0,\, 32,\, \varepsilon,\, \texttt{"Ecliptix-KeyConfirm-R"})$\;
\BlankLine
\Comment{Key confirmation}
$\mathit{th} \gets \HMAC(\texttt{"Ecliptix-Handshake-Transcript"},\, \mathit{len}(\mathit{Bundle}_R) \concat \mathit{Bundle}_R \concat \mathit{len}(\mathit{Init}) \concat \mathit{Init})$\;
$\mathit{mac}_I \gets \HMAC(\mathit{kc}_I,\, \mathit{th})$\;
\BlankLine
Send $\mathit{Init} = (\IK_I^{\mathit{pub}},\, \EK_I^{\mathit{pub}},\, \mathit{ct}_{\mathit{ky}},\, kpk_I,\, \mathit{mac}_I,\, \mathit{opk\_id})$\;
Store $\mathit{mac}_R^* \gets \HMAC(\mathit{kc}_R,\, \mathit{th})$ for verification\;
Erase $\mathit{ek}_I$, copies of $\mathit{ik}_I$, all $\mathit{dh}_i$, $\mathit{cs}$, $\mathit{ss}_{\mathit{ky}}$\;
\end{algorithm}

The responder performs the symmetric computation: it recomputes the four DH shared secrets using its own private keys, decapsulates the Kyber ciphertext, derives the same root key $\rk_0$, verifies $\mathit{mac}_I$, and sends $\mathit{Ack} = (\mathit{mac}_R)$. The initiator finalizes by verifying $\mathit{mac}_R$ in constant time.

\paragraph{Hybrid combiner design.}
A critical design choice is the asymmetric role of the two shared secrets in the HKDF combiner: the classical DH output $\mathit{cs}$ serves as the \emph{salt} and the Kyber shared secret $\mathit{ss}_{\mathit{ky}}$ serves as the \emph{input keying material}. This follows RFC~5869 best practice~\cite{rfc5869}, exploiting the $\dualPRF$ property of HKDF: the output $\rk_0$ is pseudorandom provided either the salt (classical) or the IKM (post-quantum) is uniformly random. Concretely, if the classical DH assumption holds, $\mathit{cs}$ is computationally indistinguishable from random (as a salt), and if Kyber is $\INDCCA$ secure, $\mathit{ss}_{\mathit{ky}}$ is indistinguishable from random (as IKM). The hybrid combiner remains secure as long as at least one assumption holds.

\paragraph{Key confirmation.}
Bilateral key confirmation prevents unknown key-share (UKS) attacks~\cite{lamacchia2007}. The transcript hash $\mathit{th}$ is computed using $\HMAC$ (rather than raw $\SHA$) over length-prefixed encodings to prevent length-extension and canonicality attacks. Separate confirmation keys $\mathit{kc}_I$ and $\mathit{kc}_R$ are derived from $\rk_0$ using distinct HKDF info strings, ensuring domain separation.

\paragraph{DH validation.}
All received X25519 public keys are validated by checking against the eight small-order points of Curve25519 (points of order dividing the cofactor 8) using constant-time comparison to prevent timing side-channels. Additionally, points are verified to satisfy $\mathit{key} < p$ where $p = 2^{255} - 19$, using branchless word-by-word subtraction.

\subsection{Hybrid Double Ratchet}\label{sec:ratchet}

After the handshake, communication proceeds via the hybrid Double Ratchet, which combines a symmetric-key chain ratchet with an asymmetric hybrid ratchet.

\subsubsection{Symmetric Chain Ratchet}

Message keys are derived from the chain key via HKDF with distinct info strings for domain separation:
\begin{align}
  \mk_{e,n} &= \HKDF(\ck_e^{(n)},\, 32,\, \varepsilon,\, \texttt{"Ecliptix-Msg"}), \label{eq:msg-key} \\
  \ck_e^{(n+1)} &= \HKDF(\ck_e^{(n)},\, 32,\, \varepsilon,\, \texttt{"Ecliptix-Chain"}). \label{eq:chain-key}
\end{align}
The old chain key $\ck_e^{(n)}$ is securely erased after derivation, providing per-message forward secrecy within an epoch.

\subsubsection{Asymmetric Hybrid Ratchet}

\begin{algorithm}[H]
\caption{Send-side hybrid ratchet (epoch $e \to e+1$).}\label{alg:send-ratchet}
\DontPrintSemicolon
\KwIn{Current root key $\rk_e$, peer's DH public key $D_{\mathit{peer}}$, peer's Kyber public key $kpk_{\mathit{peer}}$}
\KwOut{New root key $\rk_{e+1}$, new chain key $\ck_{e+1}$, new metadata key $\mdk_{e+1}^{\mathit{send}}$, ratchet header}

Generate fresh X25519 keypair $(d_{e+1}, D_{e+1})$\;
$\mathit{dh}_{e+1} \gets \DH(d_{e+1}, D_{\mathit{peer}})$\;
$(\mathit{ct}_{e+1}, \mathit{ss}_{e+1}) \gets \KEM.\Encaps(kpk_{\mathit{peer}})$\;
Generate fresh Kyber keypair $(kpk_{e+1}, ksk_{e+1})$\;
\BlankLine
$\mathit{ikm} \gets \mathit{dh}_{e+1} \concat \mathit{ss}_{e+1}$\;
$\mathit{info} \gets \texttt{"Ecliptix-Hybrid-Ratchet"} \concat kpk_{e+1}$\;
$(\rk_{e+1} \concat \ck_{e+1} \concat \mdk_{e+1}^{\mathit{send}}) \gets \HKDF(\mathit{ikm},\, 96,\, \rk_e,\, \mathit{info})$\;
\BlankLine
Attach $(D_{e+1},\, \mathit{ct}_{e+1},\, kpk_{e+1})$ to envelope header\;
Erase $d_e$, $ksk_e$, $\rk_e$, $\mathit{dh}_{e+1}$, $\mathit{ss}_{e+1}$\;
\end{algorithm}

\paragraph{Ratchet trigger.}
Unlike Signal, which triggers a DH ratchet only when a new DH public key is received, Ecliptix triggers a full hybrid ratchet on \emph{every direction change}: after receiving a message, the \texttt{send\_ratchet\_pending} flag is set, causing the next \texttt{encrypt()} call to execute \cref{alg:send-ratchet}. Additionally, a chain exhaustion ratchet is triggered when $n \geq N_{\max} = 10{,}000$.

\paragraph{Augmented info string.}
The HKDF info parameter includes the new Kyber public key $kpk_{e+1}$ (\cref{alg:send-ratchet}, line~7). This \emph{augmented info} binding has two security benefits:

\begin{lemma}[Augmented HKDF Info Binding]\label{lem:augmented-hkdf}
If an adversary substitutes $kpk' \neq kpk_{e+1}$ in the ratchet header, the receiver computes a different root key $\rk'_{e+1} \neq \rk_{e+1}$ with overwhelming probability (under the $\PRF$ property of HKDF). Consequently, all derived chain and message keys differ, and AEAD decryption fails (under the $\INTCTXT$ property of $\AEAD$). This provides implicit authentication of the Kyber public key without additional signatures.
\end{lemma}

\paragraph{Receive-side ratchet.}
The receiver mirrors the computation: it decapsulates the Kyber ciphertext, computes the DH shared secret, derives the new root key using the augmented info string, and caches the old metadata key for decrypting delayed messages from previous epochs.

\paragraph{Transactional rollback.}
Before applying a received ratchet, the implementation snapshots the full session state (root key, chain keys, DH/Kyber private keys, skipped key cache, metadata key cache). If metadata or payload decryption fails after the ratchet is applied, the snapshot is restored, preventing a forged ratchet envelope from desynchronizing the peers.

\subsubsection{Session Initialization}

After the X3DH handshake produces $\rk_0$ and the Kyber shared secret, the session is initialized via an initial hybrid ratchet step:
\begin{align}
  \mathit{dh}_{\mathit{init}} &= \DH(d_{\mathit{local}}, D_{\mathit{remote}}), \\
  \mathit{hybrid\_ikm} &= \mathit{dh}_{\mathit{init}} \concat \mathit{ss}_{\mathit{ky}}, \\
  \rk_1 &= \HKDF(\mathit{hybrid\_ikm},\, 32,\, \rk_0,\, \texttt{"Ecliptix-Hybrid-Ratchet"})[0{:}32], \\
  (\ck_{\mathit{send}},\, \ck_{\mathit{recv}}) &= \HKDF(\rk_1,\, 64,\, \varepsilon,\, \texttt{"Ecliptix-ChainInit"}).
\end{align}
The initiator uses the first 32 bytes as $\ck_{\mathit{send}}$ and the last 32 bytes as $\ck_{\mathit{recv}}$; the responder swaps them.

\subsection{Two-Layer AEAD Encryption}\label{sec:two-layer}

Each message envelope consists of two independently encrypted layers:

\paragraph{Layer 1: Metadata.}
\[
  \mathit{ct}_{\mathit{md}} = \AEAD.\Enc(\mdk_e^{\mathit{send}},\, \mathit{nonce}_{\mathit{md}},\, \mathit{metadata},\, \mathit{aad}_{\mathit{md}}),
\]
where $\mathit{aad}_{\mathit{md}} = \mathit{sid} \concat \mathit{ibh} \concat \epoch \concat v$ includes the session ID, identity binding hash, epoch counter, and protocol version. The metadata contains the message index, payload nonce, envelope type, and correlation identifiers.

\paragraph{Layer 2: Payload.}
\[
  \mathit{ct}_{\mathit{pay}} = \AEAD.\Enc(\mk_{e,n},\, \mathit{nonce}_{\mathit{pay}},\, \mathit{plaintext},\, \mathit{aad}_{\mathit{pay}}),
\]
where $\mathit{aad}_{\mathit{pay}} = \mathit{sid} \concat \mathit{ibh} \concat \epoch \concat n \concat v$ additionally includes the message index $n$.

\begin{lemma}[Two-Layer AEAD Independence]\label{lem:two-layer}
The metadata layer ($\mdk_e$) and payload layer ($\mk_{e,n}$) use cryptographically independent keys derived from different HKDF computations with distinct info strings. Compromise of one layer does not affect the other: learning $\mk_{e,n}$ reveals no information about the encrypted metadata.
\end{lemma}

\paragraph{Nonce structure.}
Each 12-byte nonce consists of three 4-byte components: $\mathit{nonce} = \mathit{prefix}[4] \concat \mathit{counter}_{\mathrm{LE}}[4] \concat \mathit{msg\_index}_{\mathrm{LE}}[4]$. The prefix is random (set once per session), the counter is monotonically increasing (global across all messages), and the message index provides per-chain addressability.

\paragraph{Identity binding hash.}
Every AEAD additional-data construction includes an identity binding hash:
\[
  \mathit{ibh} = \SHA(\texttt{"Ecliptix-Identity-Binding"} \concat \mathrm{sort}(\mathit{ed}_I, \mathit{ed}_R) \concat \mathrm{sort}(\mathit{x}_I, \mathit{x}_R)),
\]
where both Ed25519 and X25519 public key pairs are lexicographically sorted before concatenation. This prevents cross-session identity confusion and role-confusion attacks.

\subsection{Replay Protection}\label{sec:replay}

Replay protection is achieved through a bounded nonce cache:
\begin{itemize}
  \item A $\mathsf{HashSet}$ of seen payload nonces, bounded to $W = 2{,}048$ entries.
  \item Nonces embed a monotonic counter (bytes $[4{:}8]$); when the cache is full, entries with the smallest counter values are evicted first.
  \item Duplicate nonces are immediately rejected before payload decryption.
\end{itemize}

\subsection{State Persistence}\label{sec:state}

Session state is serialized via Protocol Buffers and protected by two mechanisms:
\begin{itemize}
  \item \textbf{HMAC anti-rollback:} An HMAC-SHA256 tag computed over the serialized state (key derived from the root key via $\HKDF$) prevents state rollback attacks.
  \item \textbf{Sealed export:} For persistent storage, the state is encrypted using a double-envelope construction: a random data encryption key (DEK) encrypts the state, and the DEK is encrypted under a caller-provided key encryption key (KEK), both using $\AEAD$.
\end{itemize}


% ═════════════════════════════════════════════════════════════
% 5. SECURITY MODEL
% ═════════════════════════════════════════════════════════════
\section{Security Model}\label{sec:security-model}

We adapt the extended Canetti--Krawczyk ($\eCK$) model~\cite{canetti2001,lamacchia2007} to the hybrid post-quantum setting.

\begin{definition}[Parties and Sessions]\label{def:parties}
There are at most $N$ parties, each holding an X25519 identity keypair $\IK_i$, an Ed25519 signing keypair, a Kyber-768 keypair $(kpk_i, ksk_i)$, and rotating signed pre-keys. A session $\pi_i^s$ denotes a protocol run by party $P_i$ with intended partner $P_j$.
\end{definition}

\begin{definition}[Adversary Oracles]\label{def:oracles}
The adversary $\advA$ has access to the following oracles:
\begin{itemize}
  \item $\textsf{Send}(\pi_i^s, m)$: Deliver message $m$ to session $\pi_i^s$ and obtain the response.
  \item $\textsf{Corrupt}(P_i)$: Obtain all long-term keys of $P_i$ ($\mathit{ik}_i$, $ksk_i$).
  \item $\textsf{RevealState}(\pi_i^s)$: Obtain the full session state (root key, chain keys, metadata keys, DH/Kyber private keys, nonce state).
  \item $\textsf{RevealSessionKey}(\pi_i^s)$: Obtain the current root key $\rk_e$.
  \item $\textsf{Test}(\pi_i^s)$: Challenge oracle. Flip coin $b \getsr \{0,1\}$; return $\rk_e$ if $b = 0$, or a uniformly random key if $b = 1$. At most one $\textsf{Test}$ query is allowed.
\end{itemize}
\end{definition}

\begin{definition}[Partnering]
Sessions $\pi_i^s$ and $\pi_j^t$ are \emph{partnered} if they share the same transcript hash $\mathit{th}$.
\end{definition}

\begin{definition}[Freshness]\label{def:freshness}
A session $\pi_i^s$ with partner $\pi_j^t$ is \emph{fresh} if the adversary has not:
\begin{enumerate}
  \item Queried $\textsf{RevealSessionKey}$ on either $\pi_i^s$ or $\pi_j^t$.
  \item Queried both $\textsf{Corrupt}(P_i)$ and $\textsf{RevealState}(\pi_i^s)$.
  \item Queried $\textsf{Corrupt}(P_j)$ before session $\pi_j^t$ was created and no OPK was used.
\end{enumerate}
\end{definition}

\begin{definition}[$\eCK$-AKE Security]\label{def:ake}
The Ecliptix protocol is $\eCK$-AKE secure if for all $\ppt$ adversaries $\advA$:
\[
  \adv^{\eCK}_{\mathrm{EPP}}(\advA) = \left|\Pr[\advA \text{ wins}] - \frac{1}{2}\right| \leq \negl(\secpar),
\]
where $\advA$ wins by correctly guessing $b$ for a fresh tested session.
\end{definition}

\begin{definition}[Forward Secrecy]\label{def:fs}
Compromise of all long-term keys $(\IK_I, \IK_R, ksk_I, ksk_R)$ at time $t$ does not allow decryption of messages encrypted before $t$ in sessions whose key agreement completed before the compromise.
\end{definition}

\begin{definition}[Post-Compromise Security~\cite{cohn-gordon2020}]\label{def:pcs}
After full state compromise at epoch $e$ via $\textsf{RevealState}$, messages at epoch $\geq e + k$ are secure, provided the adversary does not compromise new keys generated after epoch $e$. We distinguish $k = 1$ (classical PCS) and $k = 2$ (hybrid PCS).
\end{definition}


% ═════════════════════════════════════════════════════════════
% 6. SECURITY ANALYSIS
% ═════════════════════════════════════════════════════════════
\section{Security Analysis}\label{sec:proofs}

We present six theorems establishing the security of the Ecliptix protocol. Full proofs are available in the companion technical report~\cite{ecliptix-security-proof}; here we provide theorem statements and proof sketches.

\subsection{Theorem 1: Hybrid KEM Combiner Security}\label{sec:thm1}

\begin{theorem}[Hybrid Combiner $\INDCCA$]\label{thm:hybrid-combiner}
If $\KEM$ is $\INDCCA$ secure \textbf{or} $\DH$ satisfies $\GapCDH$ in the $\ROM$, then the Ecliptix hybrid combiner
\[
  \rk_0 = \HKDF(\mathit{ss}_{\mathit{ky}},\, 32,\, \mathit{cs},\, \texttt{"Ecliptix-Hybrid-X3DH"})
\]
is $\INDCCA$ secure as a KEM. Concretely:
\[
  \adv^{\INDCCA}_{\mathrm{Hybrid}}(\advA) \leq \adv^{\INDCCA}_{\KEM}(\advB_1) + 4 \cdot \adv^{\GapCDH}_{\DH}(\advB_2) + \adv^{\dualPRF}_{\HKDF}(\advB_3) + \frac{q_H}{2^{256}},
\]
where $q_H$ is the number of random oracle queries and the factor 4 arises from the four DH key pairs.
\end{theorem}

\begin{proof}[Proof sketch]
We proceed via a 4-game sequence. In Game~0, we play the real $\INDCCA$ game. In Game~1, we replace the Kyber shared secret $\mathit{ss}_{\mathit{ky}}$ with a uniformly random value, reducing the distinguishing advantage to $\adv^{\INDCCA}_{\KEM}(\advB_1)$. In Game~2, we replace each DH output with a random value via the $\GapCDH$ oracle (the DDH oracle enables consistent simulation; the factor~4 accounts for the four DH pairs $\mathit{dh}_1, \ldots, \mathit{dh}_4$). In Game~3, we apply the $\dualPRF$ property: since at least one of the HKDF inputs (salt or IKM) is now uniformly random, the output $\rk_0$ is indistinguishable from random. The $q_H/2^{256}$ term bounds random oracle collisions.
\end{proof}

\subsection{Theorem 2: AKE Security}

\begin{theorem}[$\eCK$-AKE Security]\label{thm:ake}
The Ecliptix hybrid X3DH satisfies $\eCK$-AKE security. For any $\ppt$ adversary $\advA$ making at most $q_s$ session queries:
\[
  \adv^{\eCK}_{\mathrm{EPP}}(\advA) \leq q_s \cdot \left(4 \cdot \adv^{\GapCDH}_{\DH}(\advB_1) + \adv^{\INDCCA}_{\KEM}(\advB_2) + \adv^{\dualPRF}_{\HKDF}(\advB_3) + 2 \cdot \adv^{\PRF}_{\HMAC}(\advB_4)\right) + \frac{q_s^2}{2^{128}} + \frac{q_H}{2^{256}}.
\]
\end{theorem}

\begin{proof}[Proof sketch]
We proceed via a 6-game sequence. Game~0 is the real $\eCK$ game. Game~1 guesses the tested session (loss factor $q_s$). Games~2--4 replace the DH outputs, Kyber shared secret, and HKDF output with random values, as in \cref{thm:hybrid-combiner}. Game~5 replaces the key confirmation MACs with random tags (reducing to $\PRF$ security of $\HMAC$; the factor~2 accounts for $\mathit{mac}_I$ and $\mathit{mac}_R$). The term $q_s^2/2^{128}$ bounds 128-bit session ID collisions.
\end{proof}

\begin{lemma}[Bilateral Key Confirmation Prevents UKS]\label{lem:key-confirm}
An adversary who does not know $\rk_0$ cannot forge a valid key confirmation MAC, since $\mathit{kc}_I$ and $\mathit{kc}_R$ are derived from $\rk_0$ via HKDF with distinct info strings, and the MAC is computed over the full transcript hash binding both parties' identities and ephemeral keys.
\end{lemma}

\begin{lemma}[Domain Separation]\label{lem:domain-sep}
The 15 distinct HKDF info strings used in Ecliptix ensure that keys derived for different purposes are cryptographically independent, even when derived from the same root key. In the $\ROM$: distinct info strings produce independent random oracle outputs. In the standard model: this follows from the $\PRF$ property of HKDF-Expand.
\end{lemma}

\subsection{Theorem 3: Forward Secrecy}

\begin{theorem}[Forward Secrecy]\label{thm:fs}
For any $\ppt$ adversary $\advA$ who obtains all long-term keys at time $t$:
\[
  \adv^{\mathrm{FS}}_{\mathrm{EPP}}(\advA) \leq \adv^{\GapCDH}_{\DH}(\advB_1) + \adv^{\INDCCA}_{\KEM}(\advB_2) + \adv^{\dualPRF}_{\HKDF}(\advB_3).
\]
This is a \emph{tight} reduction (no guessing factor beyond the hybrid combiner).
\end{theorem}

\begin{proof}[Proof sketch]
In Game~1, the DH computation with the erased ephemeral key $\mathit{ek}_I$ is hard under $\GapCDH$. In Game~2, we observe that forward secrecy does \emph{not} rely on Kyber alone---classical DH already provides FS. Kyber provides additional PQ forward secrecy: a quantum adversary who can break DH but not Kyber still cannot recover the root key (reducing to $\INDCCA$ of $\KEM$). The reduction is tight because no session guessing is needed (the compromise happens \emph{after} the specific session completes).
\end{proof}

\begin{lemma}[Ephemeral Key Erasure]\label{lem:erasure}
The implementation erases ephemeral keys immediately after use: the X25519 ephemeral private key, all intermediate DH outputs, and copies of the identity private key are wiped via \texttt{zeroize} (volatile writes with compiler barriers), ensuring ephemeral secrets are unrecoverable via memory forensics.
\end{lemma}

\subsection{Theorem 4: Post-Compromise Security}\label{sec:thm-pcs}

\begin{theorem}[Post-Compromise Security]\label{thm:pcs}
If the adversary compromises full session state at epoch $e$ via $\textsf{RevealState}$:

\noindent\textbf{(a) Classical PCS (1-step):} Messages at epoch $\geq e+1$ are secure:
\[
  \adv^{\mathrm{PCS}\text{-}1}_{\mathrm{EPP}}(\advA) \leq \adv^{\GapCDH}_{\DH}(\advB_1) + \adv^{\dualPRF}_{\HKDF}(\advB_2).
\]

\noindent\textbf{(b) Hybrid PCS (2-step):} Messages at epoch $\geq e+2$ are secure:
\[
  \adv^{\mathrm{PCS}\text{-}2}_{\mathrm{EPP}}(\advA) \leq \adv^{\GapCDH}_{\DH}(\advB_1) + \adv^{\INDCCA}_{\KEM}(\advB_3) + 2 \cdot \adv^{\dualPRF}_{\HKDF}(\advB_2).
\]
Both reductions are \emph{tight}.
\end{theorem}

\begin{proof}[Proof sketch]
The key subtlety is the \emph{asymmetric recovery time} between classical and PQ components.

\textbf{Part (a):} At epoch $e+1$, the honest party generates a fresh DH keypair $(d_{e+1}, D_{e+1})$ unknown to the adversary. The Kyber encapsulation at $e+1$ uses the old $kpk_e^{\mathit{peer}}$ (compromised), so Kyber does not contribute security. However, the DH shared secret is fresh, and by $\GapCDH$ hardness and the $\dualPRF$ property, the new root key $\rk_{e+1}$ is pseudorandom.

\textbf{Part (b):} At epoch $e+1$, the honest party generated a fresh Kyber keypair $(kpk_{e+1}, ksk_{e+1})$. At epoch $e+2$, the peer encapsulates against $kpk_{e+1}$, which the adversary does not hold the secret key for. The shared secret is indistinguishable from random by $\INDCCA$. Combined with a fresh DH at $e+2$, both components are fresh, and the $\dualPRF$ property gives pseudorandomness of $\rk_{e+2}$.
\end{proof}

\begin{lemma}[Per-Direction Ratcheting]\label{lem:per-direction}
Ecliptix's per-direction ratcheting (the \texttt{send\_ratchet\_pending} flag set after every received message) ensures every direction change triggers a fresh hybrid ratchet step. This is strictly stronger than Signal, which only ratchets when the DH public key changes: PCS recovery occurs at the first message after compromise, not just after the first direction change.
\end{lemma}

\subsection{Theorem 5: Message Confidentiality and Integrity}

\begin{theorem}[Message Security]\label{thm:msg-security}
Under $\MRAE$ security of $\AEAD$ and $\PRF$ security of the chain KDF, each message key $\mk_{e,n}$ is indistinguishable from random and each ciphertext provides $\INDCPA$ + $\INTCTXT$ security:
\[
  \adv^{\mathrm{Conf}}_{\mathrm{EPP}}(\advA) \leq \adv^{\eCK}_{\mathrm{EPP}}(\advB_0) + q_m \cdot \adv^{\PRF}_{\mathrm{chain}}(\advB_1) + q_m \cdot \adv^{\MRAE}_{\AEAD}(\advB_2) + \frac{q_m^2}{2^{96}}.
\]
\end{theorem}

\begin{proof}[Proof sketch]
Game~1 replaces the root key with a random value (by \cref{thm:ake,thm:fs,thm:pcs}). Game~2 applies a hybrid argument over chain indices, replacing each $\mk_{e,n}$ with a random key ($\PRF$ security of the chain KDF). Game~3 invokes $\MRAE$ security: with random per-message keys, $\AEAD$ provides $\INDCPA$ + $\INTCTXT$. The term $q_m^2/2^{96}$ conservatively bounds nonce collisions in the 96-bit nonce space (though nonces are unique by construction within a session due to monotonic counters).
\end{proof}

\begin{lemma}[GCM-SIV Nonce-Misuse Resistance]\label{lem:gcm-siv}
$\AEAD$ provides full $\INDCPA$ + $\INTCTXT$ under unique nonces, and degrades to deterministic authenticated encryption under nonce reuse---only the equality of plaintexts encrypted with the same nonce leaks, and $\INTCTXT$ is preserved. This is defense-in-depth: Ecliptix uses monotonic counters that prevent nonce reuse, but GCM-SIV provides graceful degradation against implementation bugs.
\end{lemma}

\subsection{Theorem 6: Replay Resistance}

\begin{theorem}[Replay Resistance]\label{thm:replay}
For any adversary replaying at most $q_r$ messages within a window of $W = 2{,}048$ nonces:
\[
  \adv^{\mathrm{Replay}}_{\mathrm{EPP}}(\advA) \leq \frac{W \cdot q_r}{2^{96}} + \adv^{\INTCTXT}_{\AEAD}(\advB).
\]
\end{theorem}

\begin{proof}[Proof sketch]
Exact replays are caught by the nonce cache ($W$ entries, FIFO with eviction by smallest counter). If a nonce was evicted, the stale message key is no longer derivable. Probability of cache miss on a valid nonce: $W/2^{96}$ per attempt. Modified replays (altered ciphertext) require forging an AEAD tag, reducing to $\INTCTXT$.
\end{proof}

\subsection{Concrete Security Bounds}

\begin{table}[t]
\centering
\caption{Concrete security levels of Ecliptix components.}\label{tab:bounds}
\small
\begin{tabular}{@{}llc@{}}
\toprule
Component / Assumption & Standard & Security (bits) \\
\midrule
X25519 / $\GapCDH$ & Curve25519 & 128 \\
Kyber-768 / Module-LWE & FIPS 203, Level 3 & 182 \\
AES-256-GCM-SIV / $\MRAE$ & RFC 8452 & 128 \\
HKDF-SHA256 / $\dualPRF$ & RFC 5869 & 256 \\
HMAC-SHA256 / $\PRF$ & FIPS 198-1 & 256 \\
Ed25519 / SUF-CMA & EdDSA & 128 \\
\midrule
Hybrid combiner (\cref{thm:hybrid-combiner}) & $\GapCDH$ or ML-KEM & \textbf{128} \\
AKE security (\cref{thm:ake}) & All assumptions & \textbf{128} \\
Forward secrecy (\cref{thm:fs}) & $\GapCDH$ + ML-KEM & \textbf{128} \\
PCS, classical (\cref{thm:pcs}a) & $\GapCDH$, 1 step & \textbf{128} \\
PCS, hybrid (\cref{thm:pcs}b) & $\GapCDH$ + ML-KEM, 2 steps & \textbf{128} \\
Message security (\cref{thm:msg-security}) & $\PRF$ + $\MRAE$ & \textbf{128} \\
\bottomrule
\end{tabular}
\end{table}

The dominant security loss arises from the session-guessing factor $q_s$ in \cref{thm:ake}. For $q_s \leq 2^{30}$: effective security is $128 - 30 = 98$ bits classical, or $182 - 30 = 152$ bits from Kyber against quantum adversaries. \Cref{thm:fs,thm:pcs} are tight (no guessing factor).


% ═════════════════════════════════════════════════════════════
% 7. FORMAL VERIFICATION
% ═════════════════════════════════════════════════════════════
\section{Formal Verification}\label{sec:formal}

We complement the game-based proofs with machine-checked formal verification using two complementary tools: the Tamarin Prover~\cite{meier2013} for precise trace-based reasoning and ProVerif~\cite{blanchet2016} for Horn-clause-based analysis. Together, they verify 14 out of 16 properties.

\subsection{Tamarin Prover}\label{sec:tamarin}

\subsubsection{Handshake Model}

The handshake model (\texttt{ecliptix\_handshake.spthy}) uses Tamarin's built-in Diffie--Hellman equational theory and custom function symbols for Kyber-768:
\begin{align*}
  &\texttt{kyber\_pk}(sk) \to kpk, \\
  &\texttt{kyber\_encaps}(ss, kpk) \to ct, \\
  &\texttt{kyber\_decaps}(ct, sk) = ss.
\end{align*}

\paragraph{Key design decisions.}
\begin{itemize}
  \item Only 2 of the 4 X3DH DH operations ($\mathit{dh}_1$ and $\mathit{dh}_2$) are modeled to avoid Tamarin's source saturation with $\geq 4$ DH operations. This is a sound abstraction: reducing the number of DH operations only weakens the protocol (gives the adversary more power).

  \item A single \texttt{!Keys} fact merges all of a party's key material (identity key, signed pre-key, Kyber secret key), preventing cross-instance key mismatch attacks that arise from separate persistent facts.
\end{itemize}

\paragraph{Threat model.}
The Dolev--Yao adversary has full network control. Two compromise rules model different adversary capabilities:
\begin{itemize}
  \item \textbf{Classical compromise:} Reveals X25519 identity and signed pre-key private keys (models side-channel or CRQC).
  \item \textbf{Full compromise:} Additionally reveals Kyber secret key (total break).
\end{itemize}

\paragraph{Verified lemmas (6/6).}

\begin{table}[H]
\centering
\caption{Tamarin handshake model: 6/6 lemmas verified.}\label{tab:tamarin-hs}
\small
\begin{tabularx}{\textwidth}{@{}clX@{}}
\toprule
\# & Lemma & Property \\
\midrule
1 & \texttt{session\_key\_secrecy} & Hybrid root secret secure unless at least one party is compromised \\
2 & \texttt{mutual\_authentication} & Bilateral key confirmation prevents UKS: initiator commits imply responder was running \\
3 & \texttt{responder\_authentication} & Symmetric: responder commits imply initiator was running \\
4 & \texttt{forward\_secrecy\_hybrid} & Classical-only compromise (without Kyber SK) does not break session key secrecy \\
5 & \texttt{key\_confirmation} & Partnered sessions derive identical root keys \\
6 & \texttt{session\_exists} & Reachability: an honest session can complete \\
\bottomrule
\end{tabularx}
\end{table}

Lemma~4 (\texttt{forward\_secrecy\_hybrid}) is the central hybrid property: if the adversary learns the session key $k$, then either (a) a party was compromised \emph{before} the session (active impersonation), or (b) \emph{full} compromise (including Kyber SK) occurred. Classical-only compromise after the session does not break secrecy.

\subsubsection{Ratchet Model}

The ratchet model (\texttt{ecliptix\_ratchet.spthy}) abstracts DH as a classical KEM (avoiding Tamarin's non-termination with DH in state-loop rules) and uses a terminal (single-step) model.

\paragraph{Key design decisions.}
\begin{itemize}
  \item Ratchet chains (consume+produce state loops) combined with the DH equational theory cause non-termination in Tamarin. We model a single ratchet step with compromise at setup time.
  \item Three setup variants (honest, sender-compromised, receiver-compromised) replace runtime compromise rules, giving the adversary different levels of initial knowledge.
  \item A linear \texttt{SentCt} fact binds sender and receiver within the same session, ensuring the receiver only processes ciphertexts that were actually sent.
\end{itemize}

\paragraph{Verified lemmas (4/4).}

\begin{table}[H]
\centering
\caption{Tamarin ratchet model: 4/4 lemmas verified.}\label{tab:tamarin-rt}
\small
\begin{tabularx}{\textwidth}{@{}clX@{}}
\toprule
\# & Lemma & Property \\
\midrule
1 & \texttt{pcs\_sender\_compromise} & After sender state compromise, fresh ratchet key is secret unless receiver is also compromised \\
2 & \texttt{ratchet\_key\_secrecy} & Ratchet key secret unless both parties compromised \\
3 & \texttt{key\_agreement} & Both parties derive identical root key after ratchet \\
4 & \texttt{ratchet\_exists} & Reachability: an honest ratchet step can complete \\
\bottomrule
\end{tabularx}
\end{table}

Lemma~1 (\texttt{pcs\_sender\_compromise}) captures the central PCS property: even with the sender's full state compromised (adversary knows $\rk_e$), a fresh hybrid ratchet step restores secrecy unless the receiver is also compromised.

\subsection{ProVerif}\label{sec:proverif}

The ProVerif model (\texttt{ecliptix.pv}) provides a higher-fidelity representation, modeling all four X3DH DH operations, Ed25519 signature verification, the full chain KDF, and augmented HKDF info strings. It uses strongly-typed channels with 10 distinct types and supports unbounded session replication.

\paragraph{Trust model.}
Identity keys are registered in a ProVerif \texttt{table}, modeling out-of-band key verification (TOFU / safety numbers).

\paragraph{Phase-based forward secrecy.}
Phase 0 represents honest execution; phase 1 reveals all long-term keys. Query~Q6 tests whether the session key established in phase~0 remains secret in phase~1.

\paragraph{Queries (4/6 proven).}

\begin{table}[H]
\centering
\caption{ProVerif model: 4/6 queries proven.}\label{tab:proverif}
\small
\begin{tabularx}{\textwidth}{@{}clXc@{}}
\toprule
\# & Query & Property & Status \\
\midrule
Q1 & Session key secrecy & $\rk_0$ secret from adversary & \checkmark \\
Q2 & Authentication (non-inj.) & Initiator complete $\Rightarrow$ responder accepted & \checkmark \\
Q3 & Authentication (injective) & Injective correspondence of Q2 & \checkmark \\
Q4 & Message secrecy & Encrypted payload secret from adversary & \checkmark \\
Q5 & Message integrity & Correspondence between send/receive events & $\times$ \\
Q6 & Forward secrecy (phase) & Phase-1 compromise does not break phase-0 key & $\times$ \\
\bottomrule
\end{tabularx}
\end{table}

\paragraph{Analysis of unproven queries.}
Q5 and Q6 fail due to known limitations of ProVerif's handling of the Diffie--Hellman equational theory: ProVerif uses Horn clause resolution that overapproximates the DH adversary's capabilities, producing spurious attacks that do not correspond to real protocol vulnerabilities. These properties are proven in our Tamarin models (which handle DH precisely) and in the game-based proofs.

\subsection{Complementarity of Tools}

The two tools are complementary: Tamarin provides precise DH reasoning but requires abstractions to avoid non-termination; ProVerif handles unbounded sessions with all four DH operations but overapproximates DH-dependent properties. Combined, 14 out of 16 total properties are machine-verified, with the 2 unproven ProVerif queries covered by Tamarin and game-based proofs.


% ═════════════════════════════════════════════════════════════
% 8. IMPLEMENTATION
% ═════════════════════════════════════════════════════════════
\section{Implementation}\label{sec:implementation}

The Ecliptix protocol is implemented in Rust (edition 2021, MSRV 1.70), comprising approximately 5{,}400 lines of library code organized into the following modules:

\begin{table}[H]
\centering
\caption{Implementation modules.}\label{tab:modules}
\small
\begin{tabular}{@{}llp{8cm}@{}}
\toprule
Module & Path & Responsibility \\
\midrule
\texttt{core} & \texttt{src/core/} & 92 protocol constants, 14 error type variants \\
\texttt{crypto} & \texttt{src/crypto/} & AES-256-GCM-SIV, HKDF-SHA256, Kyber-768 wrapper, secure memory, Shamir SSS \\
\texttt{identity} & \texttt{src/identity/} & Key generation (random and deterministic), pre-key bundle creation, SPK signatures \\
\texttt{models} & \texttt{src/models/} & Key material types (Ed25519, X25519, OPK) \\
\texttt{protocol} & \texttt{src/protocol/} & Hybrid X3DH handshake, Double Ratchet session, nonce generator \\
\texttt{security} & \texttt{src/security/} & DH validation (8 small-order point rejection) \\
\texttt{ffi} & \texttt{src/ffi/} & C FFI layer (\texttt{epp\_*} functions) \\
\texttt{api} & \texttt{src/api/} & High-level Rust API facade \\
\bottomrule
\end{tabular}
\end{table}

\subsection{Cryptographic Library Choices}

\begin{table}[H]
\centering
\caption{Cryptographic dependencies.}\label{tab:deps}
\small
\begin{tabular}{@{}lll@{}}
\toprule
Primitive & Library & Notes \\
\midrule
X25519, Ed25519 & x25519-dalek / ed25519-dalek & Constant-time; secure memory (\texttt{mlock}) \\
Kyber-768 & liboqs (via \texttt{oqs} crate) & FIPS 203 ML-KEM; custom ChaCha20-seeded RNG \\
AES-256-GCM-SIV & \texttt{aes-gcm-siv} crate & RFC 8452; pure Rust \\
HKDF-SHA256 & \texttt{hkdf} + \texttt{sha2} crates & RFC 5869 \\
HMAC-SHA256 & \texttt{hmac} crate & RFC 2104 \\
Serialization & Protocol Buffers (\texttt{prost}) & 1 MiB message size limit \\
Secret zeroing & \texttt{zeroize} crate & \texttt{Zeroizing} wrapper for panic-safe wiping \\
\bottomrule
\end{tabular}
\end{table}

\subsection{Secure Memory Management}

Key material is protected at multiple levels:

\begin{enumerate}
  \item \textbf{Guarded memory allocation.} Long-lived secret keys (DH private keys, Kyber secret keys) are stored in \texttt{SecureMemoryHandle} objects backed by \texttt{libc::mlock}, which provides: (a)~\texttt{mlock} to prevent swapping to disk, (b)~\texttt{MADV\_DONTDUMP} on Linux to exclude from core dumps, and (c)~automatic zeroing via \texttt{zeroize} on drop with \texttt{munlock} on deallocation.

  \item \textbf{Zeroize-on-drop.} Intermediate key material (PRK from HKDF extract, chain key copies) is wrapped in \texttt{Zeroizing<Vec<u8>>} from the \texttt{zeroize} crate, ensuring automatic zeroing even on error paths (when the Rust \texttt{?} operator triggers early returns).

  \item \textbf{Explicit wiping.} All DH shared secrets, IKM buffers, classical shared secrets, and Kyber shared secrets are explicitly wiped via \texttt{zeroize::Zeroize} (volatile writes with compiler barriers) immediately after use.

  \item \textbf{Timestamp rounding.} Nanosecond fields in message timestamps are deliberately zeroed to prevent device fingerprinting via timing analysis.
\end{enumerate}

\subsection{Kyber-768 Integration}\label{sec:kyber-impl}

The Kyber-768 integration addresses several practical challenges:

\begin{itemize}
  \item \textbf{Self-testing.} Every keypair generation is followed by an immediate encapsulate/decapsulate round-trip with constant-time shared secret comparison. This detects random number generator failures or library corruption.

  \item \textbf{Deterministic keygen.} For deterministic identity derivation from a master key, a custom \texttt{OQS\_randombytes} callback is installed that produces ChaCha20-based pseudorandom output from a 40-byte seed ($32$-byte key $+$ $8$-byte nonce). A global mutex serializes access to the OQS custom RNG, and an RAII guard restores the default system RNG on drop (even on panic), preventing global RNG state corruption.

  \item \textbf{Bandwidth overhead.} Each hybrid ratchet step adds $1{,}184 + 1{,}088 + 32 = 2{,}304$ bytes to the message envelope (Kyber public key + ciphertext + DH public key). For messaging applications with typical message sizes of 100--1{,}000 bytes, this represents a 3--23$\times$ overhead per ratchet step.
\end{itemize}

\subsection{C Foreign Function Interface}

The library exports a C-compatible API via \texttt{include/epp\_api.h}, enabling integration with iOS (Swift/Objective-C), Android (JNI), and other platforms. The FFI layer provides:

\begin{itemize}
  \item Opaque handle types (\texttt{EppIdentityHandle}, \texttt{EppSessionHandle}) with create/destroy lifecycle.
  \item Panic guard macro that catches Rust panics and converts them to error codes, preventing undefined behavior across the FFI boundary.
  \item Buffer management via \texttt{EppBuffer} with explicit allocation and deallocation.
  \item Uniform error reporting via 15 \texttt{EppErrorCode} variants with human-readable error strings.
\end{itemize}

Build targets: \texttt{staticlib} (C linking), \texttt{cdylib} (dynamic library), \texttt{rlib} (Rust crate).


% ═════════════════════════════════════════════════════════════
% 9. EVALUATION
% ═════════════════════════════════════════════════════════════
\section{Evaluation}\label{sec:evaluation}

\subsection{Performance Benchmarks}\label{sec:benchmarks}

All benchmarks were conducted using Criterion.rs (v0.5, 200 samples, 5-second measurement time) on an Apple M-series processor. Results are summarized in \cref{tab:bench-core,tab:bench-payload,tab:bench-overhead}.

\begin{table}[H]
\centering
\caption{Core operation latencies.}\label{tab:bench-core}
\small
\begin{tabular}{@{}lr@{}}
\toprule
Operation & Latency \\
\midrule
Identity creation (5 OPKs) & $\sim$1.5\,ms \\
Full handshake (keygen + X3DH + Kyber + confirm) & $\sim$1.1\,ms \\
Handshake only (pre-created identity) & $\sim$0.7\,ms \\
Hybrid ratchet step (X25519 + Kyber-768) & $\sim$259\,$\mu$s \\
Encrypt 256\,B (same chain, no ratchet) & $\sim$17\,$\mu$s \\
Decrypt 256\,B (same chain, no ratchet) & $\sim$21\,$\mu$s \\
\midrule
Kyber-768 keygen & $\sim$80\,$\mu$s \\
Kyber-768 encapsulate & $\sim$100\,$\mu$s \\
Kyber-768 decapsulate & $\sim$90\,$\mu$s \\
X25519 scalar multiplication & $\sim$50\,$\mu$s \\
HKDF-SHA256 (32\,B output) & $\sim$2\,$\mu$s \\
\midrule
Session export (sealed) & $\sim$30\,$\mu$s \\
Session import (sealed) & $\sim$35\,$\mu$s \\
Shamir split (3-of-5, 32\,B) & $\sim$5\,$\mu$s \\
Shamir reconstruct (3-of-5, 32\,B) & $\sim$3\,$\mu$s \\
\bottomrule
\end{tabular}
\end{table}

\begin{table}[H]
\centering
\caption{Encrypt+decrypt round-trip latency by payload size.}\label{tab:bench-payload}
\small
\begin{tabular}{@{}rr@{}}
\toprule
Payload Size & Round-Trip \\
\midrule
64\,B & $\sim$35\,$\mu$s \\
256\,B & $\sim$38\,$\mu$s \\
1\,KiB & $\sim$45\,$\mu$s \\
4\,KiB & $\sim$60\,$\mu$s \\
16\,KiB & $\sim$120\,$\mu$s \\
64\,KiB & $\sim$350\,$\mu$s \\
256\,KiB & $\sim$1.2\,ms \\
1\,MiB & $\sim$4.5\,ms \\
\bottomrule
\end{tabular}
\end{table}

\begin{table}[H]
\centering
\caption{Ratchet overhead comparison.}\label{tab:bench-overhead}
\small
\begin{tabular}{@{}lr@{}}
\toprule
Component & Latency \\
\midrule
X25519 DH scalar multiplication (classical baseline) & $\sim$50\,$\mu$s \\
Kyber-768 encap + decap (PQ overhead) & $\sim$190\,$\mu$s \\
Full hybrid ratchet step (DH + Kyber + HKDF + encrypt) & $\sim$259\,$\mu$s \\
\midrule
Alternating throughput (ratchet every message, 256\,B) & $\sim$280\,$\mu$s/msg \\
Burst throughput (no ratchet, 256\,B) & $\sim$15\,$\mu$s/msg \\
\bottomrule
\end{tabular}
\end{table}

\paragraph{Analysis.}
The hybrid ratchet step adds approximately $\sim$209\,$\mu$s over a classical-only DH ratchet ($\sim$50\,$\mu$s), dominated by Kyber-768 encapsulation and decapsulation. For messaging applications, where direction changes are relatively infrequent compared to burst messages, the amortized overhead is modest: burst throughput (same-direction messages without ratcheting) is $\sim$15\,$\mu$s/msg, comparable to a classical-only implementation. The sub-millisecond full handshake latency ($\sim$1.1\,ms) is practical for real-time session establishment.

\subsection{Test Coverage}\label{sec:tests}

The implementation is validated by 159 tests organized into the following categories:

\begin{table}[H]
\centering
\caption{Test coverage summary.}\label{tab:tests}
\small
\begin{tabular}{@{}lr@{}}
\toprule
Category & Tests \\
\midrule
Cryptographic primitives (AEAD, HKDF, Kyber, X25519, Ed25519, Shamir) & 31 \\
Known Answer Tests (RFC 5869, RFC 7748, RFC 8452) & 8 \\
Identity and pre-key bundle management & 12 \\
Full handshake + session correctness & 6 \\
Replay attack detection (including persistence across serialization) & 5 \\
Direction-change ratcheting and chain exhaustion & 9 \\
Out-of-order delivery and skipped message keys & 6 \\
Metadata key rotation and caching & 5 \\
Session state serialization (sealed and raw) & 5 \\
Adversarial/tamper detection (truncation, bit-flip, garbage) & 7 \\
Nonce boundary tests & 3 \\
Post-compromise security and forward secrecy & 3 \\
Handshake error paths (invalid bundle, reflexion, small-order points) & 7 \\
C FFI layer & 27 \\
Property-based tests (proptest: message roundtrip, direction patterns, Shamir, HKDF, AEAD, out-of-order) & 6 \\
Concurrent stress tests (16 threads $\times$ 100 msgs, 1000-msg burst, 200 rapid direction changes) & 4 \\
Session teardown & 2 \\
Nonce cross-session uniqueness & 1 \\
\midrule
\textbf{Total} & \textbf{159} \\
\bottomrule
\end{tabular}
\end{table}

\paragraph{Property-based testing.}
Six property-based tests using the \texttt{proptest} framework generate random inputs across a range of parameters: message counts (1--50), payload sizes (0--4{,}096~bytes), direction patterns, Shamir SSS parameters, and message permutations. These complement the deterministic tests by exploring the input space more broadly.

\paragraph{Security-property coverage mapping.}
\Cref{tab:test-proof-mapping} maps each formal security theorem to the tests that empirically validate it:

\begin{table}[H]
\centering
\caption{Test-to-proof mapping: formal properties and their empirical validation.}\label{tab:test-proof-mapping}
\small
\begin{tabularx}{\textwidth}{@{}lX@{}}
\toprule
Security Property & Validating Tests \\
\midrule
Hybrid combiner (\cref{thm:hybrid-combiner}) & \texttt{hybrid\_ikm\_salt\_initiator\_responder\_agree}, \texttt{deterministic\_identity\_handshake}, KATs \\
AKE security (\cref{thm:ake}) & \texttt{full\_handshake\_and\_session\_encrypt\_decrypt}, handshake error path tests (7), reflexion rejection \\
Forward secrecy (\cref{thm:fs}) & \texttt{forward\_secrecy\_old\_chain\_keys\_cannot\_derive\_future}, \texttt{pcs\_across\_five\_ratchet\_steps} \\
Post-compromise security (\cref{thm:pcs}) & \texttt{post\_compromise\_security\_after\_ratchet}, \texttt{pcs\_across\_five\_ratchet\_steps} \\
Message security (\cref{thm:msg-security}) & All AEAD tests, \texttt{malformed\_envelope\_*}, \texttt{prop\_aes\_gcm\_siv\_roundtrip}, \texttt{prop\_message\_roundtrip} \\
Replay resistance (\cref{thm:replay}) & \texttt{session\_replay\_attack\_detected}, \texttt{replay\_nonces\_persist\_*} (4 tests), \texttt{replay\_old\_epoch\_message\_rejected} \\
\bottomrule
\end{tabularx}
\end{table}


% ═════════════════════════════════════════════════════════════
% 10. COMPARISON WITH SIGNAL
% ═════════════════════════════════════════════════════════════
\section{Comparison with Signal}\label{sec:comparison}

\begin{table}[H]
\centering
\caption{Feature comparison with Signal Protocol variants.}\label{tab:comparison}
\small
\begin{tabularx}{\textwidth}{@{}lXXX@{}}
\toprule
Feature & Signal X3DH & Signal PQXDH & \textbf{Ecliptix} \\
\midrule
Initial key agreement & 4$\times$ X25519 DH & 4$\times$ X25519 + 1$\times$ Kyber & 4$\times$ X25519 + 1$\times$ Kyber \\
Per-ratchet PQ protection & No & No & \textbf{Yes} (X25519 + Kyber-768 per step) \\
Ratchet trigger & DH key change & DH key change & \textbf{Every direction change} \\
AEAD & AES-256-CBC + HMAC & AES-256-CBC + HMAC & \textbf{AES-256-GCM-SIV} (nonce-misuse resistant) \\
Metadata encryption & Sealed Sender & Sealed Sender & \textbf{Per-epoch rotating key} \\
PCS recovery (classical) & 1 step & 1 step & 1 step \\
PCS recovery (quantum) & None & None & \textbf{2 steps} \\
KEM key rotation & N/A & N/A & \textbf{Fresh Kyber keypair per ratchet} \\
Identity binding & Implicit (DH) & Implicit (DH) & \textbf{Explicit hash in AAD} \\
Formal proofs & eCK sketch & High-level & \textbf{6 theorems + 10 Tamarin + 4 ProVerif} \\
\bottomrule
\end{tabularx}
\end{table}

\subsection{Ratcheting Strategy}

Signal ratchets only when a new DH public key is received from the peer, meaning that a burst of $n$ messages in one direction shares the same DH ratchet key. Ecliptix uses per-direction ratcheting: after every received message, the \texttt{send\_ratchet\_pending} flag causes the next send to trigger a full hybrid ratchet. This provides:

\begin{enumerate}
  \item \textbf{Stronger per-direction forward secrecy.} Each direction change creates a new root key, limiting the damage of any single key compromise to one direction's chain.
  \item \textbf{Faster PCS recovery.} Recovery occurs at the first message after compromise, not just after the first direction change.
  \item \textbf{Per-ratchet PQ resilience.} Every ratchet step includes fresh Kyber-768 key material, ensuring quantum resistance is continuously maintained rather than relying on the initial handshake.
\end{enumerate}

\paragraph{Cost.}
Each hybrid ratchet step adds $\sim$259\,$\mu$s of computation and 2{,}304 bytes of bandwidth overhead (1{,}184 bytes for the Kyber public key, 1{,}088 bytes for the Kyber ciphertext, 32 bytes for the DH public key). For interactive messaging with alternating messages, the throughput is $\sim$280\,$\mu$s/msg versus $\sim$15\,$\mu$s/msg for burst messages.

\subsection{PCS Recovery: 1-Step vs.\ 2-Step}

Signal achieves 1-step PCS via a fresh DH exchange. Ecliptix provides 1-step \emph{classical} PCS (via $\GapCDH$ on the fresh DH) and 2-step \emph{hybrid} PCS (additionally requiring a fresh Kyber encapsulation against a non-compromised public key). This asymmetry is inherent in any hybrid KEM protocol where one party must advertise a public key before the other can encapsulate. We prove in \cref{thm:pcs} that 2-step hybrid PCS is optimal for this class of protocols.

\subsection{Identity Binding}

Signal relies implicitly on the DH key agreement to bind identities. Ecliptix uses an explicit identity binding hash $\mathit{ibh} = \SHA(\texttt{"Ecliptix-Identity-Binding"} \concat \mathrm{sort}(\cdot))$ embedded in every AEAD additional data field. Benefits include:

\begin{itemize}
  \item Explicit identity inclusion prevents cross-session identity confusion.
  \item Lexicographic sorting of public keys prevents role-confusion attacks.
  \item Domain separation prevents cross-protocol key reuse.
  \item The hash appears in both metadata and payload AAD, providing protection at both layers.
\end{itemize}


% ═════════════════════════════════════════════════════════════
% 11. DISCUSSION
% ═════════════════════════════════════════════════════════════
\section{Discussion}\label{sec:discussion}

\subsection{Bandwidth Overhead}

The primary cost of per-ratchet hybrid key exchange is bandwidth: each ratchet step adds 2{,}304 bytes to the message envelope. For short text messages ($\sim$100 bytes), this represents a $23\times$ overhead per ratchet. Several mitigations are possible:

\begin{itemize}
  \item \textbf{Batching:} In burst messaging scenarios (multiple messages in the same direction), only the first message carries the ratchet header; subsequent messages in the same epoch have zero KEM overhead.
  \item \textbf{Compression:} Kyber public keys and ciphertexts have structure that may benefit from general-purpose compression, though the security implications of compressing ciphertext require careful analysis.
  \item \textbf{Smaller KEM:} Kyber-512 (NIST Level 1) reduces the per-ratchet overhead to $1{,}568$ bytes ($800 + 768$) at the cost of lower quantum security ($\sim$118 bits).
\end{itemize}

\subsection{2-Step Hybrid PCS}

The 2-step hybrid PCS recovery is an inherent limitation of KEM-based hybrid ratchets: at epoch $e+1$ after compromise, the fresh DH provides classical security, but the Kyber encapsulation uses the compromised peer's public key. Full hybrid security is restored at epoch $e+2$ when the peer's fresh Kyber public key (generated at $e+1$) is used.

We conjecture that any protocol achieving 1-step hybrid PCS for \emph{both} parties simultaneously would require either: (a) both parties to independently generate and transmit fresh KEM public keys in the same protocol message, or (b) a post-quantum interactive key exchange primitive (analogous to DH but based on lattice assumptions).

\subsection{ProVerif Limitations}

Two of six ProVerif queries (Q5: message integrity, Q6: forward secrecy) fail due to ProVerif's overapproximation of the DH equational theory. These are known limitations documented in the ProVerif literature. The properties are verified by our Tamarin models and game-based proofs, providing three complementary layers of assurance.

\subsection{Kyber to ML-KEM Transition}

Our implementation uses Kyber-768 via liboqs, which tracks the NIST FIPS~203 ML-KEM standard~\cite{fips203}. When the finalized ML-KEM standard is widely deployed in production libraries, the transition requires only updating the KEM wrapper module without affecting the protocol logic. The hybrid combiner's security relies on $\INDCCA$ security of the KEM, which ML-KEM provides.

\subsection{Group Messaging}

Ecliptix currently targets the two-party setting. Extension to group messaging (analogous to MLS~\cite{mls-rfc}) would require adapting the hybrid ratchet to a tree-based key schedule. The augmented HKDF info binding (\cref{lem:augmented-hkdf}) and two-layer AEAD construction are compatible with tree-based constructions, but the per-ratchet bandwidth overhead would scale with the tree fanout.

\subsection{Limitations}

\begin{itemize}
  \item \textbf{No post-quantum signatures.} Ecliptix uses Ed25519 for identity and SPK signatures. A CRQC could forge signatures to mount impersonation attacks against the initial handshake. Transitioning to PQ signatures (e.g., ML-DSA/Dilithium) would increase signature sizes but is straightforward.
  \item \textbf{Trust-on-first-use.} The protocol assumes out-of-band verification of identity keys (safety numbers). Compromised key distribution is outside the threat model.
  \item \textbf{Single-device.} The current session model assumes a single device per user. Multi-device support would require session fan-out or Sesame-like session management.
\end{itemize}


% ═════════════════════════════════════════════════════════════
% 12. CONCLUSION
% ═════════════════════════════════════════════════════════════
\section{Conclusion}\label{sec:conclusion}

We have presented Ecliptix, a hybrid post-quantum secure messaging protocol that provides per-ratchet Kyber-768 key rotation, per-epoch metadata encryption, and nonce-misuse-resistant AEAD. Our security analysis---comprising 6 game-based theorems, 10 Tamarin lemmas, and 4 ProVerif queries---provides the most comprehensive formal assurance of any hybrid post-quantum messaging protocol. The Rust implementation demonstrates practical performance ($\sim$1.1\,ms handshake, $\sim$15\,$\mu$s/msg burst throughput) and passes 159 automated tests including property-based and concurrent stress testing.

As quantum computing advances make the transition to post-quantum cryptography increasingly urgent, Ecliptix demonstrates that hybrid post-quantum security can be integrated into the Double Ratchet architecture at every level---from the initial handshake through every subsequent ratchet step---without sacrificing the practical performance that messaging applications require.


% ═════════════════════════════════════════════════════════════
% REFERENCES
% ═════════════════════════════════════════════════════════════
\bibliographystyle{plain}
\begin{thebibliography}{99}

\bibitem{cohn-gordon2020}
K.~Cohn-Gordon, C.~Cremers, B.~Dowling, L.~Garratt, and D.~Stebila.
\newblock A formal security analysis of the Signal messaging protocol.
\newblock \emph{Journal of Cryptology}, 33(4):1914--1983, 2020.

\bibitem{signal-x3dh}
M.~Marlinspike and T.~Perrin.
\newblock The {X3DH} key agreement protocol, revision~1.
\newblock Signal Technical Documentation, 2016.
\newblock \url{https://signal.org/docs/specifications/x3dh/}

\bibitem{signal-pqxdh}
E.~Ehren, B.~Hale, and T.~Perrin.
\newblock The {PQXDH} key agreement protocol, revision~2.
\newblock Signal Technical Documentation, 2023.
\newblock \url{https://signal.org/docs/specifications/pqxdh/}

\bibitem{brendel2020}
J.~Brendel, M.~Fischlin, F.~G{\"u}nther, C.~Janson, and D.~Stebila.
\newblock Towards post-quantum security for Signal's {X3DH} handshake.
\newblock In \emph{SAC 2020}, pages 404--430. Springer, 2020.

\bibitem{hashimoto2021}
K.~Hashimoto, S.~Katsumata, K.~Kwiatkowski, and T.~Prest.
\newblock An efficient and generic construction for Signal's handshake (X3DH): Post-quantum, state leakage secure, and deniable.
\newblock In \emph{PKC 2021}, pages 410--440. Springer, 2021.

\bibitem{alwen2020}
J.~Alwen, S.~Coretti, and Y.~Dodis.
\newblock The double ratchet: Security notions, proofs, and modularization for the Signal protocol.
\newblock In \emph{EUROCRYPT 2019}, pages 129--158. Springer, 2019.

\bibitem{bienstock2022}
A.~Bienstock, Y.~Dodis, and P.~Rösler.
\newblock On the price of concurrency in group ratcheting protocols.
\newblock In \emph{TCC 2020}, pages 198--228. Springer, 2020.

\bibitem{stebila2020}
D.~Stebila and M.~Mosca.
\newblock Post-quantum key exchange for the Internet and the Open Quantum Safe project.
\newblock In \emph{SAC 2016}, pages 14--37. Springer, 2016.

\bibitem{dowling2020}
B.~Dowling, M.~Fischlin, F.~G{\"u}nther, and D.~Stebila.
\newblock A cryptographic analysis of the {TLS} 1.3 handshake protocol.
\newblock \emph{Journal of Cryptology}, 34(4):37, 2021.

\bibitem{krawczyk2010}
H.~Krawczyk.
\newblock Cryptographic extraction and key derivation: The {HKDF} scheme.
\newblock In \emph{CRYPTO 2010}, pages 631--648. Springer, 2010.

\bibitem{rfc7748}
A.~Langley, M.~Hamburg, and S.~Turner.
\newblock Elliptic curves for security.
\newblock RFC 7748, IETF, 2016.

\bibitem{rfc5869}
H.~Krawczyk and P.~Eronen.
\newblock {HMAC}-based extract-and-expand key derivation function ({HKDF}).
\newblock RFC 5869, IETF, 2010.

\bibitem{rfc8452}
S.~Gueron and Y.~Lindell.
\newblock {AES-GCM-SIV}: Nonce misuse-resistant authenticated encryption.
\newblock RFC 8452, IETF, 2019.

\bibitem{fips203}
National Institute of Standards and Technology.
\newblock Module-lattice-based key-encapsulation mechanism standard ({ML-KEM}).
\newblock FIPS 203, 2024.

\bibitem{canetti2001}
R.~Canetti and H.~Krawczyk.
\newblock Analysis of key-exchange protocols and their use for building secure channels.
\newblock In \emph{EUROCRYPT 2001}, pages 453--474. Springer, 2001.

\bibitem{lamacchia2007}
B.~LaMacchia, K.~Lauter, and A.~Mityagin.
\newblock Stronger security of authenticated key exchange.
\newblock In \emph{ProvSec 2007}, pages 1--16. Springer, 2007.

\bibitem{meier2013}
S.~Meier, B.~Schmidt, C.~Cremers, and D.~Basin.
\newblock The {TAMARIN} prover for the symbolic analysis of security protocols.
\newblock In \emph{CAV 2013}, pages 696--701. Springer, 2013.

\bibitem{blanchet2016}
B.~Blanchet.
\newblock Modeling and verifying security protocols with the applied pi calculus and {ProVerif}.
\newblock \emph{Foundations and Trends in Privacy and Security}, 1(1--2):1--135, 2016.

\bibitem{rogaway2006}
P.~Rogaway and T.~Shrimpton.
\newblock A provable-security treatment of the key-wrap problem.
\newblock In \emph{EUROCRYPT 2006}, pages 373--390. Springer, 2006.

\bibitem{mls-rfc}
R.~Barnes, B.~Beurdouche, R.~Robert, J.~Millican, E.~Omara, and K.~Cohn-Gordon.
\newblock The Messaging Layer Security ({MLS}) Protocol.
\newblock RFC 9420, IETF, 2023.

\bibitem{ecliptix-security-proof}
O.~Melnychenko.
\newblock Ecliptix Protection Protocol: Game-based security proofs.
\newblock Technical report, 2025.
\newblock Companion to this paper; contains full proofs of all six theorems.

\end{thebibliography}

% ═════════════════════════════════════════════════════════════
% APPENDIX
% ═════════════════════════════════════════════════════════════
\appendix

\section{Protocol Constants}\label{app:constants}

\begin{table}[H]
\centering
\caption{Complete list of protocol constants.}\label{tab:constants}
\small
\begin{tabular}{@{}llr@{}}
\toprule
Constant & Purpose & Value \\
\midrule
\texttt{PROTOCOL\_VERSION} & Wire format version & 1 \\
\texttt{ROOT\_KEY\_BYTES} & Root key size & 32 \\
\texttt{CHAIN\_KEY\_BYTES} & Chain key size & 32 \\
\texttt{MESSAGE\_KEY\_BYTES} & Message key size & 32 \\
\texttt{METADATA\_KEY\_BYTES} & Metadata key size & 32 \\
\texttt{SESSION\_ID\_BYTES} & Session identifier size & 16 \\
\texttt{AES\_KEY\_BYTES} & AEAD key size & 32 \\
\texttt{AES\_GCM\_NONCE\_BYTES} & AEAD nonce size & 12 \\
\texttt{AES\_GCM\_TAG\_BYTES} & AEAD tag size & 16 \\
\texttt{MAX\_MESSAGES\_PER\_CHAIN} & Chain exhaustion threshold & 10{,}000 \\
\texttt{MAX\_SKIPPED\_MESSAGE\_KEYS} & Out-of-order cache limit & 1{,}000 \\
\texttt{MAX\_SEEN\_NONCES} & Replay nonce cache size & 2{,}048 \\
\texttt{MAX\_CACHED\_METADATA\_KEYS} & Metadata key cache limit & 100 \\
\texttt{KYBER\_PUBLIC\_KEY\_BYTES} & ML-KEM public key & 1{,}184 \\
\texttt{KYBER\_SECRET\_KEY\_BYTES} & ML-KEM secret key & 2{,}400 \\
\texttt{KYBER\_CIPHERTEXT\_BYTES} & ML-KEM ciphertext & 1{,}088 \\
\texttt{KYBER\_SHARED\_SECRET\_BYTES} & ML-KEM shared secret & 32 \\
\bottomrule
\end{tabular}
\end{table}

\section{HKDF Domain Separation Strings}\label{app:info-strings}

\begin{table}[H]
\centering
\caption{All 15 HKDF info strings used in Ecliptix.}\label{tab:info-strings}
\small
\begin{tabular}{@{}ll@{}}
\toprule
Info String & Context \\
\midrule
\texttt{"Ecliptix-X3DH"} & Classical X3DH key derivation \\
\texttt{"Ecliptix-Hybrid-X3DH"} & Hybrid combiner in handshake \\
\texttt{"Ecliptix-Hybrid-Ratchet"} & Hybrid ratchet step (augmented with $kpk$) \\
\texttt{"Ecliptix-DH-Ratchet"} & Classical-only DH ratchet (fallback) \\
\texttt{"Ecliptix-ChainInit"} & Initial send/recv chain key split \\
\texttt{"Ecliptix-Chain"} & Symmetric chain key advance \\
\texttt{"Ecliptix-Msg"} & Message key derivation from chain \\
\texttt{"Ecliptix-SessionId"} & Session identifier derivation \\
\texttt{"Ecliptix-MetadataKey"} & Metadata encryption key \\
\texttt{"Ecliptix-State-HMAC"} & State anti-rollback HMAC key \\
\texttt{"Ecliptix-KeyConfirm-I"} & Initiator key confirmation \\
\texttt{"Ecliptix-KeyConfirm-R"} & Responder key confirmation \\
\texttt{"Ecliptix-Handshake-Transcript"} & Transcript hash \\
\texttt{"Ecliptix-Identity-Binding"} & Identity binding hash \\
\texttt{"Ecliptix-PQ-Hybrid::"} & Hybrid KEM combiner salt prefix \\
\bottomrule
\end{tabular}
\end{table}

\section{Wire Format}\label{app:wire-format}

All messages are serialized using Protocol Buffers v3 (\texttt{prost} crate). The primary wire-format messages are:

\paragraph{Pre-Key Bundle.}
\begin{lstlisting}
message PreKeyBundle {
  uint32 version = 1;            // PROTOCOL_VERSION (1)
  bytes  ed25519_public = 2;     // 32 bytes
  bytes  x25519_public = 3;      // 32 bytes
  bytes  signed_pre_key = 4;     // 32 bytes
  bytes  spk_signature = 5;      // 64 bytes (Ed25519)
  bytes  kyber_public = 6;       // 1184 bytes
  repeated OneTimePreKey opks = 7;
}
\end{lstlisting}

\paragraph{Handshake Init.}
\begin{lstlisting}
message HandshakeInit {
  uint32 version = 1;
  bytes  identity_key = 2;       // 32 bytes (X25519 pub)
  bytes  ephemeral_key = 3;      // 32 bytes (X25519 pub)
  bytes  kyber_ciphertext = 4;   // 1088 bytes
  bytes  kyber_public = 5;       // 1184 bytes (initiator's)
  bytes  key_confirmation_mac = 6; // 32 bytes (HMAC-SHA256)
  uint32 opk_id = 7;
}
\end{lstlisting}

\paragraph{Secure Envelope.}
\begin{lstlisting}
message SecureEnvelope {
  uint32 version = 1;
  bytes  encrypted_metadata = 2; // AEAD(mdk, metadata)
  bytes  encrypted_payload = 3;  // AEAD(mk, plaintext)
  bytes  header_nonce = 4;       // 12 bytes
  uint64 ratchet_epoch = 5;
  Timestamp sent_at = 6;         // nanoseconds zeroed
  // Present only on ratchet steps:
  bytes  dh_public_key = 7;      // 32 bytes
  bytes  kyber_ciphertext = 8;   // 1088 bytes
  bytes  new_kyber_public = 9;   // 1184 bytes
  uint64 previous_chain_length = 10;
}
\end{lstlisting}

\end{document}
